\section{模型假设}
\begin{enumerate}[label=\arabic*.]
\item 官方转写与视频中的话语顺序一致，转写不改变话语先后。依据是该转写由赛题统一提供并与音轨逐条对应。这一条使词级强制对齐的单调路径约束成立，未通过字符目标构造的词也可按同一顺序单独记录。
\item 同一原词的子词共享父词时间区间，帧特征在其时间支撑区间内近似为分段常量。依据是词为发音的自然单位，而音频帧率远高于词长：前者使子词继承父词的音视频锚点，后者使式\eqref{q1:eq:pool}的交叠积分退化为加权求和。
\item 附件2至附件4的对齐版本中，同一序列位置的三个模态共享官方给定的时间对应关系，局部缺失长度以锚点数度量。依据是这些附件按\texttt{aligned\_50}组织，位置本身即对应单元。问题二与问题三因此只需处理观测可用性，不必重新估计时间对应。
\item 内容位置的文本\file{[UNK]}与全零音视频向量表示观测不可用，而非有效的特征取值。依据是附件3对局部缺失的定义即为特征值全为零的连续区间。填充、观测与缺失三种状态由此可由掩码唯一确定。
\item 人工遮挡只改变输入证据，不改变样本标签，既定情感目标在证据减少时仍具有可预测性。依据是遮挡由程序在训练阶段随机生成，原始标签不受其影响。问题二据此在附件2上构造缺失条件下的监督信号。
\item 附件2训练集估计的标准化统计量适用于验证集与附件3、附件4，推理期间模型权重与标准化统计量保持固定。依据是各附件由同一特征提取流程生成，维度与量纲一致。该约定避免推理阶段引入测试数据信息，并使跨附件指标在同一尺度下可比。
\end{enumerate}
